% Generated draft: requires mathematical and source review.
\documentclass[11pt]{article}
\usepackage[margin=25mm]{geometry}
\usepackage{fontspec}
\setmainfont{Latin Modern Roman}
\usepackage{amsmath,amssymb,mathtools,unicode-math}
\setmathfont{Latin Modern Math}
\usepackage{enumitem,needspace,xurl,hyperref}
\setlength{\parindent}{0pt}
\setlength{\parskip}{5pt}
\setlength{\emergencystretch}{3em}
\newcommand{\N}{\mathbb{N}}
\newcommand{\Z}{\mathbb{Z}}
\newcommand{\Q}{\mathbb{Q}}
\newcommand{\R}{\mathbb{R}}
\newcommand{\C}{\mathbb{C}}
\newcommand{\F}{\mathbb{F}}
\newcommand{\A}{\mathbb{A}}
\newcommand{\PP}{\mathbb{P}}
\newcommand{\E}{\mathbb{E}}
\newcommand{\eps}{\varepsilon}
\newcommand{\dd}{\,\mathrm{d}}
\begin{document}
\section*{UnsolvedMath ID: 30004621}
\textbf{Source problem code:} catalogue entry
\subsection{Definitions and mathematical statement}
Let \(g\) be an integer with \(g\geq 3\). The supplied record denotes by \(\mathcal A_g\) the moduli space of principally polarized abelian varieties of dimension \(g\). It does not specify the base field, whether the moduli stack or a coarse moduli space is intended, or whether a compactification is required.

The question recorded for permanent UnsolvedMath ID \(30004621\) is whether, for every such \(g\),
\[
\operatorname{Eff}_2(\mathcal A_g)
\stackrel{?}{=}
\operatorname{Cone}\bigl(S_A,S_F,S_D,S_1,S_2\bigr).
\]
Here the right-hand side expresses the proposed generation by the five named classes: under a real-vector-space convention, it would mean
\[
\left\{a_A S_A+a_F S_F+a_D S_D+a_1 S_1+a_2 S_2:
 a_A,a_F,a_D,a_1,a_2\in\R_{\geq0}\right\}.
\]
This interpretation is conditional: the record supplies neither definitions of these five classes nor their dependence on \(g\), their ambient cycle-class space, the equivalence relation on cycles, or a convention concerning closure of the effective cone.

There is also an unresolved dimensional incompatibility. The original statement and literature assessment call the target a \emph{codimension-two} effective cone, whereas the clean statement uses \(\operatorname{Eff}_2\) and calls the proposed generators \emph{surface classes}. For the usual moduli space, \(\dim\mathcal A_g=g(g+1)/2\); hence surfaces have codimension \(g(g+1)/2-2\), which is not two for any integer \(g\geq3\). Thus two-dimensional cycles and codimension-two cycles cannot be identified here. The displayed equality preserves the supplied question symbolically; it cannot be made into a uniquely specified mathematical conjecture from this record without resolving this incompatibility and defining the five classes and their ambient space.
\subsection{Short English statement}
For every integer \(g\geq3\), the supplied question asks whether five named classes generate the indicated effective cone associated with the moduli space of principally polarized abelian varieties. The record leaves the classes undefined and conflates surfaces with codimension-two cycles, so the intended conjecture requires clarification. Its open status is supplied dated curation.
\subsection{Sources}
\begin{itemize}
\item Local registry: \texttt{data/unsolvedmath/problems.json}, permanent UnsolvedMath ID \(30004621\), source problem code \texttt{catalogue entry}, in the Oberwolfach Reports Open Problems collection.
\item Supplied source citation: \emph{Moduli spaces and Modular forms (2022)}, Oberwolfach Reports, pp.\ 301--359, DOI \texttt{10.4171/owr/2021/6}. \url{https://doi.org/10.4171/owr/2021/6}. The background describes Oberwolfach Reports as documenting mathematical workshops and their open-problem discussions. The record gives proposed year \(2021\), but names no proposer.
\item Supplied literature source: Grushevsky--Hulek, ``effective surfaces.'' \url{https://arxiv.org/abs/2007.02995}. No further bibliographic details or definitions of the five classes are supplied.
\item Supplied literature source: Oberwolfach report. \url{https://doi.org/10.4171/owr/2021/6}.
\item Literature assessment dated exactly \texttt{2026-08-22}: the supplied curation labels the problem open, describes the all-genera generation assertion as conjectural, and reports that available work proves extremality of several rays without proving the asserted equality. This is supplied dated curation, not an independently verified literature assessment; the background explicitly requests independent verification of time-sensitive status and cited claims.
\item The registry labels the statement \texttt{corrected\_verified} and reports that it was checked against the official Oberwolfach report. That verification claim is part of the supplied record and has not been independently confirmed for this draft.
\end{itemize}
\end{document}
